\usepackage[margin=25mm]{geometry}
\usepackage[T1]{fontenc}
\usepackage{lmodern}
\usepackage{microtype}
\usepackage{amsmath,amssymb,mathtools,bm}
\usepackage{braket}
\usepackage{booktabs,longtable,tabularx,array}
\usepackage{graphicx}
\usepackage{siunitx}
\usepackage[dvipsnames]{xcolor}
\usepackage{enumitem}
\usepackage{fancyhdr}
\usepackage{caption}
\usepackage{subcaption}
\usepackage{float}
\usepackage{needspace}
\usepackage{listings}
\usepackage[most]{tcolorbox}
\usepackage{tikz}
\usetikzlibrary{arrows.meta,positioning,shapes.geometric}
\usepackage[backend=biber,style=numeric-comp,sorting=nyt]{biblatex}
\usepackage{hyperref}
\usepackage[nameinlink,noabbrev]{cleveref}

\addbibresource{bibliography/references.bib}

\definecolor{SemBlue}{HTML}{123B66}
\definecolor{SemLightBlue}{HTML}{EEF5FB}
\definecolor{SemGreen}{HTML}{246B4A}
\definecolor{SemOrange}{HTML}{A44A17}
\definecolor{CodeBackground}{HTML}{F5F7FA}

\hypersetup{
  colorlinks=true,
  linkcolor=SemBlue,
  citecolor=SemGreen,
  urlcolor=SemOrange,
  pdftitle={Semiconductor Physics: Models, Computation, and Validation},
  pdfauthor={Eugene Joseph M. Ragasa}
}

\sisetup{
  scientific-notation=true,
  separate-uncertainty=true,
  per-mode=symbol
}

\setlist{itemsep=0.25em,topsep=0.4em}
\setlength{\parindent}{0pt}
\setlength{\parskip}{0.55em}
\setcounter{secnumdepth}{3}
\setcounter{tocdepth}{2}

\pagestyle{fancy}
\setlength{\headheight}{14pt}
\fancyhf{}
\fancyhead[L]{\small\nouppercase{\leftmark}}
\fancyhead[R]{\small\thepage}
\renewcommand{\headrulewidth}{0.3pt}

\newtcbtheorem[number within=chapter]{definition}{Definition}{
  enhanced,
  breakable,
  colback=SemLightBlue,
  colframe=SemBlue,
  colbacktitle=SemBlue,
  coltitle=white,
  fonttitle=\bfseries,
  boxrule=0.8pt,
  arc=1mm,
  before skip=0.8em,
  after skip=0.8em
}{def}

\newtcolorbox{modelcheck}[1][]{
  enhanced,
  breakable,
  colback=green!4,
  colframe=SemGreen,
  title={Model check},
  fonttitle=\bfseries,
  #1
}

\newtcolorbox{boundarybox}[1][]{
  enhanced,
  breakable,
  colback=orange!5,
  colframe=SemOrange,
  title={Model boundary},
  fonttitle=\bfseries,
  #1
}

\lstset{
  language=Python,
  basicstyle=\ttfamily\small,
  backgroundcolor=\color{CodeBackground},
  frame=single,
  rulecolor=\color{black!20},
  showstringspaces=false,
  columns=fullflexible,
  keepspaces=true,
  breaklines=true,
  numbers=left,
  numberstyle=\scriptsize\color{black!55},
  xleftmargin=2em,
  framexleftmargin=1.5em
}

\newcommand{\dd}{\mathop{}\!\mathrm{d}}
\newcommand{\kb}{k_{\mathrm B}}
\newcommand{\Ec}{E_{\mathrm C}}
\newcommand{\Ev}{E_{\mathrm V}}
\newcommand{\Eg}{E_{\mathrm g}}
\newcommand{\Ef}{E_{\mathrm F}}
\newcommand{\mui}{\mu_{\mathrm i}}
\newcommand{\Nc}{N_{\mathrm C}}
\newcommand{\Nv}{N_{\mathrm V}}
\newcommand{\ii}{\mathrm i}

% Every published figure identifies the student-facing notebook that reproduces it.
\newcommand{\figurenotebookmark}{\protect\footnotemark}
\newcommand{\figurenotebooktext}[1]{%
  \footnotetext{\textit{Student reproduction notebook:}
    \href{run:#1}{\nolinkurl{#1}}.}%
}
